% defined-in: spivak (page 541)
% === SEMANTIC MAPPING ===
% 1. Variables & Types: 
%    - $a_n \colon \NaturalNumbers \TermTo \RealNumbers$ (Sequence of coefficients)
%    - $f \colon \RealNumbers \TermTo \RealNumbers$ (Power series function)
%    - $x \colon \RealNumbers$ (Variable)
%    - $n, m \colon \NaturalNumbers$ (Indices)
%    - $\epsilon \colon \RealNumbers$ (Positive real)
% 2. Data Structures: 
%    - $\Series{\sum_{n=0}^{\infty} a_n x^n}$ represents the infinite series.
%    - $\ClosedInterval{0, 1}$ represents the domain of uniform convergence.
% 3. Implicit Bounds: 
%    - The series is defined for $x \in \ClosedInterval{0, 1}$.
%    - The convergence of $\sum a_n$ is assumed as the premise.
% ========================

% entity-id: prop-abels-theorem-power-series
% entity-type: prop
% macro: \AbelsTheoremPowerSeries
% args: 0

\begin{proposition}[Abel's Theorem]
\[
\TerForall a \colon \NaturalNumbers \TermTo \RealNumbers \left( \Series{\sum_{n=0}^{\infty} a_n} \text{ converges} \TermImplication 
\begin{aligned}
& \left( \Series{\sum_{n=0}^{\infty} a_n x^n} \text{ converges uniformly on } \ClosedInterval{0, 1} \right) \TermConjunction \\
& \left( \FunctionLimitAtPoint{f}{1} = \sum_{n=0}^{\infty} a_n \right) \TermConjunction \\
& \left( \Continuous{f \text{ on } \ClosedInterval{0, 1}} \right)
\end{aligned}
\right)
\]
\end{proposition}

\begin{proof}[RU]
Пусть $\Series{\sum_{n=0}^{\infty} a_n}$ сходится. Согласно критерию Коши для рядов, для любого $\epsilon > 0$ существует $N \in \mathbb{N}$ такое, что для всех $m > n > N$ выполняется $\RealAbsoluteValue{\sum_{k=n}^{m} a_k} < \epsilon$. Используя лемму Абеля, для любого $x \in \ClosedInterval{0, 1}$ получаем $\RealAbsoluteValue{\sum_{k=n}^{m} a_k x^k} \leq \RealAbsoluteValue{\sum_{k=n}^{m} a_k} < \epsilon$. Это доказывает равномерную сходимость ряда на $\ClosedInterval{0, 1}$. В силу теоремы о равномерной сходимости, предел функции $f(x)$ при $x \to 1$ равен сумме ряда $\sum_{n=0}^{\infty} a_n$, а сама функция $f$ непрерывна на $\ClosedInterval{0, 1}$.
\end{proof}

\begin{proof}[EN]
Suppose $\Series{\sum_{n=0}^{\infty} a_n}$ converges. By the Cauchy criterion for series, for any $\epsilon > 0$, there exists $N \in \mathbb{N}$ such that for all $m > n > N$, $\RealAbsoluteValue{\sum_{k=n}^{m} a_k} < \epsilon$. Applying Abel's Lemma, for any $x \in \ClosedInterval{0, 1}$, we have $\RealAbsoluteValue{\sum_{k=n}^{m} a_k x^k} \leq \RealAbsoluteValue{\sum_{k=n}^{m} a_k} < \epsilon$. This establishes the uniform convergence of the series on $\ClosedInterval{0, 1}$. By the theorem on uniform convergence, the limit of the function $f(x)$ as $x \to 1$ is equal to the sum of the series $\sum_{n=0}^{\infty} a_n$, and the function $f$ is continuous on $\ClosedInterval{0, 1}$.
\end{proof}